\section{Sensitivity to the Modeling Assumptions}
\label{sec:robustness}
% SWEVO expansion: tab:robust-final is back in the main text (it was moved to the supplement by editor cut C2 in the
% MPCE version); the full re-evaluation stays in the supplement (Table S-robust). tab:F-bench and tab:F-equiv are
% captured in 08_beyond.tex (\CaptureSuppTables{analysis/mpce_supp_direction.tex}).
The benchmark fixes a power curve, a wake model, a minimum spacing, a boundary rule and a discretization of the wind rose. We re-evaluated every feasible final layout of the main comparison, without re-optimization, under a cubic power curve~\cite{Carrillo2013} (also with a 25~m/s cut-out), the uncalibrated Gaussian wake of~\cite{Bastankhah2014} ($K_{\rm G}=0.04$), and a finer direction resolution (Table~\ref{tab:robust-final}; all methods: Table~\ref{S-tab:robust}). A re-evaluation asks whether the delivered layouts keep their order under another model; it does not show which layouts the methods would find if they optimized that model. The spacing and the boundary rule are examined on the final layouts and in a separate SSA experiment.
%% generated by mpce_results.py -- do not edit by hand
\begin{table}[!t]
\centering
\caption{Robustness to the Benchmark Model: Final Layouts Re-Evaluated, Not Re-Optimized}
\label{tab:robust-final}
\scriptsize\setlength{\tabcolsep}{2.4pt}
\begin{tabular}{lcccccc}
\toprule
& \multicolumn{2}{c}{$\Delta$ (\%)} & & & & \\
\cmidrule(lr){2-3}
Model & DS I & DS II & Rank & Best & $\bar\tau$ & Same \\
\midrule
Benchmark & +0.0 & +0.0 & 1.82 & PSO-VNS & -- & -- \\
Cubic curve & -17.7 & -36.9 & 1.82 & PSO-VNS & 0.95 & 94 \\
Cubic, 25 m/s cut-out & -21.5 & -37.0 & 1.80 & PSO-VNS & 0.95 & 93 \\
Gaussian wake & -0.3 & -0.4 & 1.97 & PSO & 0.66 & 54 \\
\bottomrule
\end{tabular}
\par\vspace{2pt}\parbox{\columnwidth}{\scriptsize All feasible final layouts of the 68 cases. $\Delta$: mean relative change of the objective (\%); Rank: average rank of PSO-VNS; Best: best-ranked method; $\bar\tau$: mean Kendall correlation between the benchmark and the alternative ordering within a case; Same: cases (\%) with unchanged best method. Ranks of all methods: supplementary material.}
\end{table}


\subsection{Power Curve}
\label{sec:powercurve}
The benchmark power curve of Kusiak and Song is replaced by a cubic curve between cut-in and rated speed~\cite{Carrillo2013}, with and without a cut-out at 25~m/s. The cubic curve lowers the expected power by \NCubicOverI\% (Data Set~I) and \NCubicOverII\% (Data Set~II), so our energy values are benchmark values, but it leaves the ranking intact ($\bar\tau=\NTauCubic$; same best method in \NSameCubic\% of the cases; PSO-VNS first, also with the cut-out). Both curves increase with the wind speed below rated, so a layout with smaller wake deficits in the frequent directions gains under either; the change of the curve shifts the level of the objective far more than the order of the layouts within a case.
% CHECK-FINAL [C21]: robustness.Cubic/CubicCutout: tau >= 0.9, PSO-VNS rank position 1

\subsection{Wake Model}
\label{sec:gaussian}
The Gaussian wake (Section~\ref{S-sec:S-gauss}) differs from the Jensen cone of the benchmark in its smooth lateral profile, in its expansion rate ($K_{\rm G}=0.04$ vs.\ $K=0.075$; the two are not calibrated to each other) and, unlike the cone test of the benchmark, gives no deficit to turbines upstream of the wake source. The mean objective changes little (Table~\ref{tab:robust-final}), but the wake loss of the same final layouts differs between the two models by \NXModelShiftMean~pp on average (median \NXModelShiftMedian~pp), \NXModelShiftRatio{} times the equivalence margin (Section~\ref{sec:setup-stats}), and by more than the margin for \NXModelShiftAboveMarginPct\% of the \NXModelShiftNLayouts{} layouts. The per-case difference PSO-VNS${}-{}$PSO changes much less, by \NXModelShiftPairMean~pp on average (\NXModelShiftPairJensen{} vs.\ \NXModelShiftPairGauss~pp). The ordering is less stable than under the cubic curve ($\bar\tau=\NTauGauss$), and PSO is narrowly first (\NFGaussFifteenRankPSO{} vs.\ \NFGaussFifteenRankPSOVNS); with 1\textdegree{} direction bins, PSO-VNS is first under the Gaussian wake as well (Section~\ref{sec:direction}). The comparison of PSO-VNS and PSO is therefore more robust to the wake model than the absolute wake losses, which depend on the model by several times the margin.
% CHECK-EXTRA [X30]: Jensen vs Gaussian wake loss of the same layouts: mean \NXModelShiftMean pp (median \NXModelShiftMedian; \NXModelShiftRatio x margin; > margin for \NXModelShiftAboveMarginPct % of \NXModelShiftNLayouts layouts); per-case PSO-VNS - PSO changes by \NXModelShiftPairMean pp on average (below the margin)
% CHECK-FINAL [C22]: robustness.Gauss (15 deg): tau < 0.8, PSOC first, PSOBV second (1.91 vs 1.97); CHECK-DIR [F09]: PSO-VNS first under the Gaussian wake at 1 deg

\subsection{Minimum Spacing}
\label{sec:spacing}
The optimized layouts often lie on the $4D$ constraint: \NSpFiveAll\% of the feasible final layouts also satisfy $5D$ and \NSpSixAll\% $6D$ (\NSpFiveLarge\% and \NSpSixLarge\% for $N=11$--15), so they do not transfer to larger spacings without re-optimization (constructive capacity of the sites: Table~\ref{S-tab:capacity}). The share falls with $N$: on sites of fixed size, more turbines leave less room beyond the minimum distance. The rankings of this paper therefore hold for $\ell_{\min}=4D$; PSO-VNS and PSO were not re-run at larger spacings, and the archived re-optimization at $5D$ and $6D$ with the original SSA and LX-SSA code (Table~\ref{S-tab:spacing-authors}) shows how the attainable objective changes, not how the methods compare.
% CHECK-FINAL [C23]: spacing_share.n11_15.pct_5d < 10 ("often lie on the 4D constraint")
% verified (manual): shares 49 % / 36 % (all N) vs 4.1 % / 0.1 % (N = 11-15) from mpce_numbers.tex (\NSp...); constructive capacity at 4D/5D/6D: Table S-capacity

\subsection{Boundary Rule}
\label{sec:boundary}
All methods clip every coordinate to the bounding square and leave the circle to the penalty (Section~\ref{sec:constraints}). The boundary rule matters: an SSA that projects turbines radially onto the circle instead of clipping the coordinates attains higher mean objectives than the original SSA in all \NBoundCases{} representative cases ($p\le\NBoundMaxP$). A radial projection returns every turbine to the site directly, whereas clipping leaves turbines in the corners of the square outside the circle, where only the penalty pushes them back. Because the rule changes the results of a method, all methods use the same rule, and the comparison is one of search strategies under a common repair, not of repair operators.
% CHECK-FINAL [C41]: boundary_rule n_higher == n_cases == 6 and max_p < 0.01

\subsection{Direction Resolution}
\label{sec:direction}
The benchmark evaluates the wake only at the centers of 24 direction bins of 15\textdegree, and a pair of turbines can be placed so that no bin center lies in its narrow wake interval (Section~\ref{sec:model}). We therefore split each 15\textdegree{} bin into 3 (5\textdegree) or \NFSubBins{} (1\textdegree) sub-bins with the bin's Weibull parameters and an equal share of its frequency, so that the wake-free power is unchanged, and re-evaluated all \NFNLayouts{} feasible final layouts (Table~\ref{tab:F-bench}).
% CHECK-DIR [F01]: one sub-bin reproduces the benchmark objective; wake-free objective identical for 1, 3 and 15 sub-bins
The 1\textdegree{} rose raises the mean wake loss of every method, by \NFLossRiseMin~pp (\NFLossRiseMinBy) to \NFLossRiseMax~pp (\NFLossRiseMaxBy; PSO-VNS: \NFLossPSOVNSFifteen\% to \NFLossPSOVNSOne\%): optimized layouts exploit the gaps between the coarse bins, those of PSO and the VNS-based methods most. Most of the rise appears already with 5\textdegree{} sub-bins. PSO-VNS still ranks first (\NFRankPSOVNSOne; PSO \NFRankPSOOne, the only method not significantly different), also under the Gaussian wake at 1\textdegree{} (\NFGaussOneRankPSOVNS{} vs.\ \NFGaussOneRankPSO); PSO's narrow lead under the Gaussian wake with 15\textdegree{} bins (\NFGaussFifteenRankPSO{} vs.\ \NFGaussFifteenRankPSOVNS; $\bar\tau=\NTauGauss$) is a property of the coarse rose. Beyond the top two the order changes ($\bar\tau=\NFTauOne$; same best method in \NFSameBestOne\% of the cases): MS-SLSQP moves from sixth to third (average rank \NFRankMSSLSQPFifteen{} to \NFRankMSSLSQPOne), and its gap to PSO-VNS shrinks from \NFSLSQPGapFifteen{} to \NFSLSQPGapOne~pp but remains significant.
% CHECK-DIR [F03]: loss rise of every method, PSO-VNS 1.322 -> 2.335; [F04]: most of the change already with 5-deg sub-bins (> 70 % for PSO-VNS); [F05]: VNS-based methods and PSO rise most, DE/MS-SLSQP least ("VNS-based methods most"); [F06]: PSO-VNS first at 5 and 1 deg, PSO the only one not significantly different (Holm); [F09]: Gaussian: PSO first at 15 deg, PSO-VNS first at 1 deg; [F07]: MS-SLSQP 6th -> 3rd, tau 0.52, same best 37 %; [F08]: MS-SLSQP gap shrinks from 0.435 to 0.135 pp, still significant (Holm)
\SuppTable{tab:F-bench}

For PSO-VNS vs.\ PSO (Table~\ref{tab:F-equiv}), the case-mean difference becomes \NFPairMeanOne~pp (90\% CI \NFPairCIOne): a small but significant advantage ($p=\NFPairPOne$), not equivalence at $\pm\NEqMargin$~pp (minimal margin \NFPairMinMarginOne~pp); the same holds with 5\textdegree{} sub-bins (\NFPairMeanFive~pp, 90\% CI \NFPairCIFive). The sign of the per-case difference changes in \NFSignChangesOne{} of the \NFSignChangesBase{} cases with a nonzero difference under both resolutions, so the per-case winner often depends on the discretization. On the \NFNonTrivN{} non-trivial cases (PSO-VNS wake loss of at least \NFNonTrivThr\%), equivalence is shown at neither resolution, and the Data Set~II $N\ge10$ advantage remains (\NFDSIINTenWinsOne{} of \NFDSIINTenN{} cases). The component-analysis conclusions hold, but the small SSA gain over RS-VNS is no longer significant (Table~\ref{S-tab:F-abl}). The order beyond the top two and the equivalence verdict are thus properties of the discretized benchmark; the first place of PSO-VNS is not.
% CHECK-DIR [F10]: PSO-VNS - PSO -0.043, CI [-0.062, -0.025], not equivalent, p < 0.05; [F11]: same with 5-deg sub-bins; [F12]: per-case sign changes in 26 of 60 cases; [F13]: non-trivial cases: equivalence at neither resolution; [F14]: DS II N >= 10 10/10; [F16]: component-analysis conclusions hold (SSA-VNS and LX-SSA-VNS worse than RSD-VNS, PSO-VNS better than RSD-VNS, RS-VNS and VNS, PSO-VNS best of nine); [F17]: SSA-VNS vs RS-VNS not significant at 1 deg
\SuppTable{tab:F-equiv}
